\section*{Chapter \ref{ch:g12:equilibrium} --- Chemical Equilibrium: Quotient and Constant}

\begin{solution}{exo:g12:equilibrium:1}
(a) $Q = \dfrac{[\ce{NH3}]^2 (c^\circ)^2}{[\ce{N2}][\ce{H2}]^3}$;
(b) $Q = \dfrac{[\ce{CH3COO-}][\ce{H3O+}]}{[\ce{CH3COOH}]\,c^\circ}$ (water,
the \omterm{def:g6:solutions-and-solubility:solute}{solvent}, is left out);
(c) $Q = [\ce{Ca^{2+}}][\ce{CO3^{2-}}]/(c^\circ)^2$ (the solid is left out);
(d) $Q = \dfrac{[\ce{FeSCN^{2+}}]\,c^\circ}{[\ce{Fe^{3+}}][\ce{SCN-}]}$.
\end{solution}

\begin{solution}{exo:g12:equilibrium:2}
$\tau = 0.020/0.050 = 0.40$: not total.
\end{solution}

\begin{solution}{exo:g12:equilibrium:3}
$Q = 2 < K$: forward. $Q = 50 > K$: reverse. $Q = K$: no change.
\end{solution}

\begin{solution}{exo:g12:equilibrium:4}
The composition is constant, but the forward and reverse reactions go on
at the same rate: \omterm{def:g7:atoms-and-molecules:molecule}{molecules} keep reacting both ways.
\end{solution}

\begin{solution}{exo:g12:equilibrium:5}
No, $K$ is unchanged (the \omterm{def:g10:reaction-progress-table:system}{final state} is the same); the time to reach
equilibrium is shorter.
\end{solution}

\begin{solution}{exo:g12:equilibrium:6}
$\dfrac{x^2}{(2-x)^2} = 4$, so $\dfrac{x}{2-x} = 2$ and
$x = \qty{1.33}{mol}$: \qty{0.67}{mol} of acid and of \omterm{def:g11:functional-groups:alcohol}{alcohol},
\qty{1.33}{mol} of \omterm{def:g11:functional-groups:carboxylic-acid}{ester} and of water (again two thirds).
\end{solution}

\begin{solution}{exo:g12:equilibrium:7}
$Q = (0.5 \times 0.5)/(0.5 \times 0.5) = 1 < 4$: forward, more \omterm{def:g11:functional-groups:carboxylic-acid}{ester}
forms.
\end{solution}

\begin{solution}{exo:g12:equilibrium:8}
About \qty{0.52}{mol} from the acid side and \qty{0.79}{mol} from the
\omterm{def:g11:functional-groups:carboxylic-acid}{ester} side. After about \qty{100}{h} both curves are on the dashed line:
equilibrium.
\end{solution}

\begin{solution}{exo:g12:equilibrium:9}
With $y$ mol hydrolysed: $\dfrac{y^2}{(1-y)^2} = \dfrac{1}{4}$, so
$\dfrac{y}{1-y} = \dfrac{1}{2}$, $y = \frac{1}{3}$: $\frac{2}{3}$ mol of
\omterm{def:g11:functional-groups:carboxylic-acid}{ester} and of water, $\frac{1}{3}$ mol of acid and of \omterm{def:g11:functional-groups:alcohol}{alcohol}. The same
final \omterm{def:g6:pure-substances-and-mixtures:mixture}{mixture} as the esterification.
\end{solution}

\begin{solution}{exo:g12:equilibrium:10}
A solid is a pure phase, whose ``concentration'' does not change: it is
left out of $Q$. Adding more solid does not change $Q$, so it does not
shift the equilibrium (as long as some solid is present).
\end{solution}

\begin{solution}{exo:g12:equilibrium:11}
$K$ fixes the ratio of carbon dioxide in the gas to that in the water.
Opened, the gas above the water escapes into the \omterm{def:g4:air-a-mixture-of-gases:air}{air}: the carbon dioxide
in the gas falls, $Q < K$, and more gas leaves the water, until almost
none is left.
\end{solution}

\begin{solution}{exo:g12:equilibrium:12}
At equilibrium: acid and \omterm{def:g11:functional-groups:alcohol}{alcohol} $\frac{1}{3}$, \omterm{def:g11:functional-groups:carboxylic-acid}{ester} and water
$\frac{2}{3}$. Removing half of the water leaves $\frac{1}{3}$:
$Q = \dfrac{(2/3)(1/3)}{(1/3)(1/3)} = 2 < 4$, forward. With $y$ more
reacting: $(2/3 + y)(1/3 + y) = 4(1/3 - y)^2$, that is
$27y^2 - 33y + 2 = 0$, $y = 0.064$: the \omterm{def:g11:functional-groups:carboxylic-acid}{ester} rises to \qty{0.73}{mol}.
\end{solution}

\begin{solution}{exo:g12:equilibrium:13}
The quotient of the reverse reaction has numerator and denominator
exchanged: $Q' = 1/Q$, so at equilibrium $K' = 1/K$. Hydrolysis:
$1/4 = 0.25$.
\end{solution}

\begin{solution}{exo:g12:equilibrium:14}
$Q = (2 \times 2)/(1 \times 1) = 4 = K$: the \omterm{def:g6:pure-substances-and-mixtures:mixture}{mixture} is already in
equilibrium; its composition does not change.
\end{solution}

\begin{solution}{exo:g12:equilibrium:15}
$n - 0.95 = 0.9025 / (4 \times 0.05) = 4.51$, so $n = \qty{5.5}{mol}$ of
\omterm{def:g11:functional-groups:alcohol}{alcohol} per \omterm{def:g10:the-mole:amount}{mole} of acid.
\end{solution}

\begin{solution}{pb:g12:equilibrium:1}
\textbf{1.} \ce{CH3-COOH + CH3-CH2-OH <=> CH3-COO-CH2-CH3 + H2O}.

\textbf{2.} It is not total: it stops in an equilibrium where both
directions go on.

\textbf{3.} $Q = \dfrac{n_{\text{ester}}\, n_{\text{water}}}{n_{\text{acid}}\, n_{\text{alcohol}}}$.

\textbf{4.} $x_{\max} = \qty{1}{mol}$; $x_f = \frac{2}{3}$ mol.

\textbf{5.} $\tau = 0.67$.

\textbf{6.} Acid and \omterm{def:g11:functional-groups:alcohol}{alcohol} $\frac{1}{3}$ mol each, \omterm{def:g11:functional-groups:carboxylic-acid}{ester} and water
$\frac{2}{3}$ mol each.

\textbf{7.} $K = \dfrac{(2/3)^2}{(1/3)^2} = 4$.

\textbf{8.} One third hydrolysed: \omterm{def:g11:functional-groups:carboxylic-acid}{ester} and water $\frac{2}{3}$, acid
and \omterm{def:g11:functional-groups:alcohol}{alcohol} $\frac{1}{3}$: the same \omterm{def:g6:pure-substances-and-mixtures:mixture}{mixture}, $Q_{eq} = 4$.

\textbf{9.} It is the same reaction at the same temperature, and $K$
depends only on the temperature.

\textbf{10.} Acid $1 - x$; \omterm{def:g11:functional-groups:alcohol}{alcohol} $3 - x$; \omterm{def:g11:functional-groups:carboxylic-acid}{ester} $x$; water $x$.

\textbf{11.} $x^2 = 4(1 - x)(3 - x) = 4(3 - 4x + x^2)$, so
$3x^2 - 16x + 12 = 0$.

\textbf{12.} $x = (16 \pm \sqrt{256 - 144})/6 = (16 \pm 10.58)/6$: 0.903
or 4.43; only $x_f = \qty{0.903}{mol}$ lies between 0 and 1.

\textbf{13.} \qty{90}{\percent}.

\textbf{14.} $M = \qty{88.0}{g/mol}$: $0.903 \times 88.0 = \qty{79.5}{g}$.

\textbf{15.} Each removal of water lowers the numerator of $Q$, which
falls below $K$: the system moves forward again.

\textbf{16.} The reaction would go on until the acid is used up:
\qty{100}{\percent}.

\textbf{17.} Ethanol is cheaper, and its excess is easily separated from
the \omterm{def:g11:functional-groups:carboxylic-acid}{ester} (and recycled).

\textbf{18.} Only the time: a reaction that gives out almost no heat has
a constant that hardly changes with temperature, and heating only makes
it reach equilibrium sooner.

\textbf{19.} About \qty{90}{\percent} of the acid.
\end{solution}
